% ============================================================
%  hw1_statement.tex
%  Exact re-typesetting of hw1.pdf (MATH 6660, Homework 1, Fall 2026).
%  Statement only -- nothing added, nothing removed.
%  Build:  pdflatex hw1_statement.tex   (run twice for the (1) references)
% ============================================================
\documentclass[11pt,leqno]{article}

% geometry chosen to reproduce the measured layout of hw1.pdf:
%   text block x: 54.0 .. 551.0 pt,  title ink top y=81.7,  folio y=723.0
\usepackage[letterpaper,left=54pt,right=61pt,top=78pt,bottom=90pt,
            headheight=30pt,headsep=14pt,footskip=30pt]{geometry}
\usepackage{amsmath,amssymb}
\usepackage{enumitem}
\usepackage{fancyhdr}
\usepackage{xcolor}
\usepackage[colorlinks=true,linkcolor={red!55!black},citecolor={red!55!black},
            urlcolor={red!55!black},pdfborder={0 0 0}]{hyperref}

\pagestyle{fancy}
\fancyhf{}
\fancyhead[R]{\textbf{Assigned:} Friday Aug. 28, 2026\\
              \textbf{Due:} Friday Sept. 11, 2026}
\fancyfoot[C]{\thepage}
\renewcommand{\headrulewidth}{0.4pt}

\newcommand{\LL}{\mathcal{L}}

\setlength{\parindent}{0pt}

% measured from hw1.pdf: "1." at x=54.0, its text at x=67.9, "(b)" at x=71.9
\setlist[enumerate,1]{leftmargin=13.9pt,labelsep=5.5pt,topsep=3pt,itemsep=4pt,parsep=0pt}
\setlist[enumerate,2]{leftmargin=24.4pt,labelsep=5.5pt,topsep=3pt,itemsep=4pt,parsep=0pt}

\begin{document}

\begin{center}
  {\large\textsc{MATH 6660: Homework 1}}\\
  {\large\textsc{Fall 2026}}
\end{center}

\begin{enumerate}

% ================================================================
\item Consider the following 2-point BVP:
\begin{gather*}
  \LL\left(u\right) := -\bigl(a(x)\,u'\bigr)' = f(x) \quad \text{on} \quad 0<x<1,\\
  u(0) = 0, \quad u(1) = 0,
\end{gather*}
where $f\in C^{0}(0,1)$, $a\in C^{1}(0,1)$, and $0<a_{1}\le a(x)\le a_{2}<\infty$.
In class, we showed that for any two continuous functions on $(0,1)$, $u$ and
$v$, that satisfy the boundary conditions, $u(0)=v(0)=u(1)=v(1)=0$, the
operator, $\LL$, satisfies Green's formula:
\[
  \int_{0}^{1}\left[u\LL\left(v\right)-v\LL\left(u\right)\right]dx = 0.
\]

\begin{enumerate}

\item Consider the associated eigenvalue problem:
\begin{equation}\label{eq:1}
  \LL\left(\phi\right) = \lambda\,\phi,
\end{equation}
where $\lambda$ is an eigenvalue and $\phi(x)$ is the corresponding
eigenfunction. Prove that $\lambda$ must be real.
\textbf{HINT:} start by showing that if $(\lambda,\phi(x))$ is an
eigenvalue/eigenvector pair, then so is
$(\overline{\lambda},\overline{\phi}(x))$, where $\overline{\,\cdot\,}$ means
complex conjugation.

\item Prove that two eigenfunctions, $\phi_{1}$ and $\phi_{2}$, corresponding to
two different eigenvalues, $\lambda_{1}$ and $\lambda_{2}$, respectively, where
$\lambda_{1}\neq\lambda_{2}$, are \textbf{orthogonal}:
\[
  \int_{0}^{1}\phi_{1}(x)\,\phi_{2}(x)\,dx = 0.
\]

\item Prove that all the eigenvalues of \eqref{eq:1} are strictly positive:
$\lambda>0$. \textbf{HINT:} Multiply \eqref{eq:1} by $\phi$ and integrate over
$(0,1)$.

\item Let $a(x)\equiv1$. Consider the special case where the load function, $f$,
is chosen as the Dirac delta centered at $s$:
\[
  \LL\left(G\right) = -G'' = \delta\left(x-s\right), \quad G(0;s)=G(1;s)=0.
\]
Compute by hand the solution to the above problem: $G(x;s)$. This is known as
the \textbf{Green's function}.

\item Let $a(x)\equiv1$. From the previous part, show that $G(x;s)=G(s;x)$.
Using an approriate application of Green's formula, show that for a general load
function, $f\in C^{0}(0,1)$, the solution to
\[
  \LL\left(u\right) = -u'' = f(x), \quad u(0)=u(1)=0
\]
is given by
\[
  u(x) = \int_{0}^{1} f(s)\,G(s;x)\,ds.
\]

\end{enumerate}

% ================================================================
\item Consider the following scalar product and the associated norm:
\[
  (v,\,w) := \int_{0}^{1} v(x)\,w(x)\,dx
  \qquad\text{and}\qquad
  \|v\| := \sqrt{(v,\,v)}.
\]
Prove Cauchy's inequality:
\[
  \left|\,(v,\,w)\,\right| \le \|v\|\|w\|.
\]
\textbf{HINT:} Construct the function: $P(t):=\|v-tw\|^{2}$, and show that $P$
is a quadratic polynomial. Think about what you can say about the roots of $P$.

\bigskip

% ================================================================
\item Consider the following boundary value problem with periodic boundary
conditions:
\begin{align*}
  \textbf{ODE} &: \quad -\, u'' + q(x)\,u = f(x),\\
  \textbf{BC1} &: \quad u(-\pi) = u(\pi),\\
  \textbf{BC2} &: \quad u'(-\pi) = u'(\pi),
\end{align*}
where $q>0$ $\forall x\in[-\pi,\pi]$.

\begin{enumerate}
\item Recast this equation into a variational problem, stating the trial and
test function spaces.

\item Interpret the variational problem as an energy minimization problem,
clearly stating the energy functional. Prove that the variational problem and
the energy minimization problems are equivalent.
\end{enumerate}

\end{enumerate}

\end{document}
